\documentclass{situationist-preprint}

\title{What It Is Like to Be a Bat: A Structural Specification}
\msid{MND-03-a}
\msfamily{MND-03}
\msdiscipline{Philosophy of Mind}
\mstype{Principal manuscript}
\mskeywords{bats; echolocation; Nagel; phenomenal structure; structuralism; qualia; objective phenomenology; spatial experience}

\begin{document}
\maketitle

\begin{abstract}
Nagel argued that we cannot know what it is like to be a bat, because the facts about bat experience are accessible only from a point of view we cannot occupy, and that this shows a limit on objective description. We argue that a large part of the answer can nonetheless be given, and we give it. Echolocating bats have a sonar-structured spatial field: an egocentric, three-dimensional array of located surfaces in which range is presented directly rather than inferred, which is sampled in discrete pulses whose rate the bat controls, whose resolution rises as the bat closes on a target, and in which, for some species, relative velocity is a dimension of the presented objects. The argument rests on a mirroring principle, on which the structure of a creature's phenomenal space mirrors the structure of the discriminations it uses to guide action, and on the account of qualia in MND-10, on which a relational description of a quality space is complete up to structure-preserving permutation and leaves out the intrinsic character of each quality. Much of what it is like to be a bat is therefore specifiable in structural terms; not all of it is. The residue is exactly what MND-10 predicts, and it is smaller than Nagel supposed.
\end{abstract}

\section{The problem}\label{s:problem}

\citet{Nagel1974} chose bats because they are close enough to us on the phylogenetic tree that few would doubt they have experience, and far enough from us in their dominant mode of perception that we have no idea what that experience is like. Most microchiropteran bats find their way and their prey by emitting high-frequency calls and analysing the returning echoes. Nagel's claim was that, although we may know everything about the physics of the calls and the neurophysiology of the bat's auditory system, we cannot know what the resulting experience is like for the bat, because that is a fact accessible only from a type of point of view, and ours is of another type. Extrapolating from our own case does not help, since our own case contains nothing like echolocation; imagining hanging upside down and eating insects tells us what it would be like for us to behave as a bat does, not what it is like for a bat to be a bat.

The question \enquote{What is it like to be a bat?} can be read in three ways. On the \emph{existential} reading it asks whether there is anything it is like, that is, whether bats are conscious. On the \emph{characterizing} reading it asks for a description of their experience that a reader who is not a bat could understand and assess. On the \emph{acquaintance} reading it asks us to know the experience as the bat does, by having it. Nagel took the first answer for granted, held that the third is unavailable to us, and argued that the second is unavailable too, at least with our present concepts, though he left open the project of an \enquote{objective phenomenology} that would describe experience in terms intelligible to beings who cannot have it.

We argue that the characterizing answer can be given to a large extent, and we give it: what it is like to be an echolocating bat is like having a sonar-structured spatial field. The specification is structural. It says how the bat's experience is organized: its dimensions, their metrics, their resolution and their temporal sampling. It does not say, and on our view cannot say, what each place in that organization is like intrinsically. The thesis therefore has two parts, and the second part is as important as the first: much of the bat's experience can be specified in structural terms, though not all.

The thesis depends on the account of qualia defended in MND-10, on which qualia exist and are intrinsic to the experiences that have them, while their structure can be described relationally. It is a test case for that account, in the sense that the account predicts both how far a structural specification can go and where it must stop. It also depends on MND-23 for the existential answer, and on MND-04 for an account of what humans lack with respect to the residue.

\section{History of the debate}\label{s:history}

\paragraph{Nagel's argument and its target.} Nagel's paper was directed against reductive physicalism, not against the existence of facts about bat experience. The argument has the following shape. Every reductive account of a phenomenon abstracts from the point of view of particular observers in order to reach a more objective description. Experience is essentially tied to a point of view. So reduction, which abstracts from points of view, will leave out what is essential to experience. The bat case serves to make vivid that the facts about experience outrun our concepts: if there are facts we cannot even formulate, there are facts no physical description we can understand will capture. Nagel's conclusion was that physicalism, though perhaps true, is a position we cannot now understand. The first-person concepts that experience requires for its full grasp were later developed into the knowledge argument of \citet{Jackson1982} and the debate over phenomenal concepts; MND-04 argues against the knowledge argument that Mary acquires abilities and concepts but no new fact.

\paragraph{Objective phenomenology.} At the end of his paper Nagel proposed a programme: to develop concepts for describing the subjective character of experience in a form comprehensible to beings incapable of having those experiences, beginning with structural features of perception. He suggested, as an illustration, that one might try to convey to a person blind from birth something of what it is like to see by describing the structure of visual space. The proposal was not developed by Nagel himself, and it has been the least discussed part of the paper. It is the part that the present manuscript carries out for the bat.

\paragraph{Deflationary replies.} \citet[ch.~14]{Dennett1991} replied that the question can in principle be answered by third-person science. If we learned in sufficient detail what the bat can and cannot discriminate, what it can and cannot use, and how its states guide its behavior, we would know what it is like to be a bat, and the residue that Nagel thinks we would miss is nothing. \citet{Hacker2002} argued that the phrase \enquote{what it is like to be} is grammatically confused when used to define consciousness: the question \enquote{what is it like?} asks for a comparison or an attitude, and in the bat case it has a dull and answerable sense in which the bat's life is like flying about in the dark catching insects. Both replies deny that there is a residue. We agree with Dennett that much can be said, and disagree that nothing remains; Hacker's reading changes the subject (Objection~\ref{o:hacker}).

\paragraph{Akins and the neuroethology of the bat.} \citet{Akins1993} took Nagel's question to the neurophysiological literature and found that, for the moustached bat in particular, much of the processing serves narrowly defined tasks: tracking the range and relative velocity of one target, with the call adjusted to keep the returning echo in a narrow frequency band. She suggested that there may be no representation of a world of objects in such a bat at all, and therefore nothing it is like to perceive one: what it is like to be a bat may be very different from having a field of surfaces. Akins's argument is the most serious challenge to the structural answer, because it accepts that structure determines character and disputes what the structure is. We take it up in Section~\ref{s:cases} and Objection~\ref{o:akins}.

\paragraph{Structuralism about qualities.} A separate tradition holds that sensory qualities are individuated by their places in a quality space, defined by relations of similarity and difference that discrimination tasks can reveal \citep{Clark1993,Rosenthal2010}. On this view a complete structural specification of the bat's quality spaces would leave nothing out, since there is nothing to a quality but its place. Recent work has developed formal models of mental qualities in this spirit, with structure carried by orderings and metrics on quality spaces \citep{Lee2021}. Structuralism agrees with Dennett that there is no residue but locates the answer in quality spaces rather than in behavior. MND-10 argues against it on the ground of Newman's objection, which is the ground on which we rest the second half of our thesis.

The positions thus form a spectrum. Nagel: the characterizing answer is not available. Dennett and the structuralists: the characterizing answer is available and complete. Akins: an answer is available, but it is not the one we expect. Our position: the characterizing answer is available, it is the expected one in its main lines, and it is incomplete in a specific and predictable way.

\section{What bats do with sound}\label{s:bio}

Echolocation in bats was established by Griffin and Galambos in the 1940s and described in detail by \citet{Griffin1958}. Most microchiropteran bats emit calls through the mouth or nose and receive the echoes through ears with large and mobile pinnae. The physics of the echo carries information of several kinds, and there is experimental and neurophysiological evidence that bats extract and use each.

\emph{Range.} The delay between call and echo gives the distance of the reflecting surface: roughly six milliseconds per metre. Bats discriminate range differences of the order of a centimetre or less, and in the auditory cortex of the moustached bat there are neurons tuned to specific call--echo delays, arranged so that delay is mapped systematically across the cortical surface \citep{Suga1990}. \emph{Direction.} Azimuth is given by differences in intensity and timing between the two ears, elevation by direction-dependent spectral filtering by the pinna and tragus. \emph{Size and texture.} Echo amplitude, corrected for range, varies with the reflecting area; the spectral fine structure of the echo varies with the depth structure of the surface, so that surfaces of different texture return echoes of different spectral composition \citep{Simmons1989}. \emph{Relative velocity.} In bats with long constant-frequency components in their calls, such as horseshoe bats and the moustached bat, the Doppler shift of the echo gives the relative velocity of the target, and the periodic modulations of the echo produced by a beating insect wing allow fluttering prey to be detected against clutter. These bats lower their call frequency as they fly so as to keep the echo within a narrow band to which a disproportionate share of their cochlea and auditory cortex is tuned, an \enquote{acoustic fovea} \citep{Suga1990}.

Two features of the bat's use of these cues are especially important for what follows. First, the bat controls its sampling. In search flight a bat may call ten times a second; on approach to prey it calls faster, and in the final \enquote{terminal buzz} it may call two hundred times a second, with shorter and broader-band calls. Sampling rate, and with it temporal resolution, rises as the target is approached. The bat also directs its sonar beam and adjusts the timing of its calls to the scene, so that echolocation is an active process in which the animal probes the environment \citep{UlanovskyMoss2008,MossSurlykke2010}. Second, the information extracted is integrated into representations of the spatial layout that guide flight. Bats avoid fine wires in the dark, navigate cluttered vegetation, and intercept prey whose trajectory changes; they also learn and remember the spatial layout of familiar places, and they combine echolocation with vision where light allows.

\section{Framework}\label{s:framework}

\begin{definition}[Phenomenal structure]\label{d:struct}
The \emph{phenomenal structure} of a modality of a creature's experience is the set of dimensions along which experiences in that modality vary, together with the relations of order, distance, betweenness and part--whole among values along those dimensions, and the way the experienced values change over time.
\end{definition}

\begin{definition}[Structural specification]\label{d:spec}
A \emph{structural specification} of a modality is a description of its phenomenal structure that mentions no quality except as the occupant of a place in that structure. Following MND-10, it is a relational description of the modality's quality space.
\end{definition}

\begin{definition}[Spatial field]\label{d:field}
An experience presents a \emph{spatial field} when its phenomenal structure includes a manifold of experienced locations, ordered by experienced relations of direction and distance from the subject, at which qualities are presented as located.
\end{definition}

\begin{definition}[Sonar-structured spatial field]\label{d:sonar}
A spatial field is \emph{sonar-structured} when (i) its locations are fixed by the reflection of the subject's own emissions, so that only surfaces that return emissions are presented at them; (ii) range from the subject is a dimension presented with high resolution, rather than reconstructed from other dimensions; (iii) the field is refreshed in discrete samples at a rate the subject controls, and its resolution varies with that rate; and (iv) the qualities presented at locations include properties of reflecting surfaces, such as extent and texture, and, where the emissions permit it, relative motion.
\end{definition}

The principles on which the argument rests are three.

\begin{principle}[Consciousness]\label{p:cons}
Bats are conscious: there is something it is like to be a bat.
\end{principle}

Principle~\ref{p:cons} is a consequence of the conclusion of MND-23 that mammals are conscious. The behavioral and neural markers on which that conclusion rests are present in bats, and the structures in question are homologous with our own.

\begin{principle}[Mirroring]\label{p:mirror}
If a creature is conscious and uses a modality to guide flexible action, then the phenomenal structure of that modality mirrors the structure of the discriminations the creature uses in that guidance: two stimuli the creature discriminates and uses as different are experienced as different along a dimension of the modality; the order and metric of the discriminations are the order and metric of the experienced values; and the integration of discriminated values into representations of a common layout is matched by their presentation in a common field.
\end{principle}

Principle~\ref{p:mirror} is the principle by which we already describe the experience of other human beings, and the principle by which psychophysics describes our own. We know that the visual field of another person is roughly elliptical, that its resolution falls off from the fovea, that it is interrupted by saccades of which the subject is not aware, and so on, from the structure of the discriminations that the person makes and uses. The principle is restricted to discriminations used in flexible guidance. Discriminations made by processes whose outputs guide only reflexes, or that are not available to the creature's control of action, need not be mirrored in its experience; blindsight, in which discriminations are made but are not used spontaneously by the subject, is the standard illustration in the human case. The restriction is also what answers the objection that the principle makes experience depend on the theorist's chosen description of the stimulus: the relevant discriminations are those the creature uses, and which those are is a question about the creature.

\begin{principle}[Intrinsicness]\label{p:intr}
Qualia are intrinsic to the experiences that have them. A relational description of a quality space is complete up to structure-preserving permutation: it fixes the places in the structure and leaves open which quality occupies each place.
\end{principle}

Principle~\ref{p:intr} is the central result of MND-10. Its first clause is the theorem of that manuscript; its second is the proposition that a relational description is complete up to automorphism, which MND-10 derives from Newman's point against Russell's structural realism \citep{Newman1928}.

\section{The argument}\label{s:argument}

\begin{lemma}\label{l:spatial}
The bat's echolocative experience presents a spatial field.
\end{lemma}

\begin{proof}
By Principle~\ref{p:cons}, the bat is conscious. The bat uses echolocation to guide flexible action: it avoids obstacles it has not met before, intercepts prey on changing trajectories and adapts its flight to clutter (Section~\ref{s:bio}). The discriminations it uses are discriminations of the direction and distance of reflecting surfaces, and they are integrated into representations of a common layout, since the bat's flight responds to the arrangement of several surfaces at once, as in passing between branches or between a wire and a wall. By Principle~\ref{p:mirror}, its echolocative experience has dimensions of direction and distance with the order and metric of these discriminations, and the discriminated values are presented in a common field. A manifold of experienced locations ordered by direction and distance from the subject, at which qualities are presented as located, is a spatial field (Definition~\ref{d:field}).
\end{proof}

\begin{lemma}\label{l:sonar}
The spatial field of Lemma~\ref{l:spatial} is sonar-structured.
\end{lemma}

\begin{proof}
We take the four clauses of Definition~\ref{d:sonar} in turn, applying Principle~\ref{p:mirror} to the discriminations described in Section~\ref{s:bio}.

(i) The discriminations are made on echoes of the bat's own calls. A surface that returns no echo, because it lies in the acoustic shadow of a nearer one or absorbs the call, is not discriminated by echolocation and is not used in its guidance. So, within the modality, only reflecting surfaces are presented at locations. The analogy is to vision, in which only surfaces that reflect or emit light are presented; the disanalogy is that the illumination is the bat's own, so that the field exists only while the bat calls and extends only as far as its calls return.

(ii) Range is discriminated directly, from echo delay, with a fine resolution, and is mapped in the auditory cortex as a dimension in its own right \citep{Suga1990}. In human vision, distance is recovered from a combination of cues, such as binocular disparity, accommodation, occlusion and texture gradients, none of which is a direct measure of it, and its resolution falls rapidly beyond a few metres. By Principle~\ref{p:mirror}, range is a dimension of the bat's field presented with the order and metric of the bat's range discriminations: fine, and not reconstructed from other experienced dimensions.

(iii) The bat samples the scene in discrete calls at a rate it controls, and the resolution of its discriminations varies with that rate, rising in the approach and terminal phases of pursuit (Section~\ref{s:bio}). The temporal structure of the discriminations is therefore that of a field refreshed in discrete samples whose rate and resolution vary with the bat's action. Whether the field is experienced as discontinuous or as continuous with discrete updates is a further question of structure, which the evidence about the bat's integration of successive echoes could in principle decide; on either answer, the refresh rate and the resolution co-vary as the definition requires.

(iv) Echo amplitude and spectral fine structure allow the discrimination of the extent and texture of surfaces, and in constant-frequency bats the Doppler shift and the modulation produced by beating wings allow the discrimination of relative velocity and of flutter. These discriminations are used in the choice and pursuit of prey. By Principle~\ref{p:mirror}, they are experienced as dimensions along which the qualities presented at locations vary. Extent, texture and, for these bats, motion are thus properties presented at the places of the field.
\end{proof}

\begin{theorem}\label{t:main}
What it is like to be an echolocating bat, in respect of its echolocative experience, is like having a sonar-structured spatial field. The structure of this field can be specified relationally; the intrinsic character of the qualities that occupy it cannot.
\end{theorem}

\begin{proof}
The first sentence is Lemma~\ref{l:sonar}. For the second: Lemmas~\ref{l:spatial} and \ref{l:sonar}, together with the further detail of Section~\ref{s:cases}, give a description of the field's dimensions, metrics, resolution and temporal sampling that mentions no quality except as the occupant of a place, which is a structural specification (Definition~\ref{d:spec}). By Principle~\ref{p:intr}, such a specification is complete up to structure-preserving permutation of the qualities in each quality space, and it leaves open which quality occupies each place. Since qualia are intrinsic, which quality occupies a place is a further fact about the bat's experience. So the intrinsic character of the qualities is not specified.
\end{proof}

The second half of Theorem~\ref{t:main} is not a confession of ignorance about bats in particular. It is the same limit that applies to any structural specification of anyone's experience, including our own, as MND-10 argues: a relational description of human color space is complete up to automorphism and does not say what red is like. What is distinctive about the bat is that we who read the specification have no experiences with which to fill the places from our own acquaintance.

\begin{proposition}\label{p:share}
The part of the bat's experience that can be specified structurally includes everything about it that can be conveyed to a subject who lacks the modality, and the residue is the intrinsic character of the occupants of its quality spaces.
\end{proposition}

\begin{proof}
What can be conveyed by description to someone who lacks a modality is what can be stated in terms that do not require acquaintance with its qualities; by MND-10, such terms describe a quality only as the occupant of a place in a structure, so what can be conveyed is structural. Conversely, everything structural can be conveyed, since a relational description can be understood by anyone who grasps the relations, as a mathematician can understand the description of a space without inhabiting it. By Theorem~\ref{t:main}, the remainder is the intrinsic character of the qualities.
\end{proof}

Proposition~\ref{p:share} vindicates Nagel's programme of objective phenomenology within limits that Nagel himself suggested: the structural features of experience are describable to those who lack it. It also vindicates the core of his scepticism, relocated. What we cannot know is not \enquote{what it is like to be a bat} in general, but which qualities occupy the places of a structure that we can describe. In the terms of MND-04, the residue is what can be grasped only under phenomenal concepts, which humans lack for echolocation as Mary in her room lacks them for color. Lacking these concepts is not lacking access to any fact that a physical description omits; it is lacking a mode of presentation of facts that are physically describable.

\begin{corollary}\label{c:comparison}
The question whether bat echolocation is more like human hearing or more like human vision has a determinate structural answer: in its spatial organization it is more like vision. Whether its qualities resemble auditory or visual qualities intrinsically is not settled by the structural answer.
\end{corollary}

\begin{proof}
Human hearing presents sound sources at coarse directions with poor range resolution; it presents sounds, which have pitch, loudness and timbre, as events at or emanating from places rather than as surfaces. Human vision presents a field of surfaces at fine directions, with properties such as extent and texture, refreshed continuously. By Lemmas~\ref{l:spatial} and \ref{l:sonar}, the bat's field presents surfaces at fine directions and ranges, with extent and texture; so its structure is closer to that of vision, with the differences listed in Definition~\ref{d:sonar}. By Principle~\ref{p:intr}, structural similarity does not fix intrinsic similarity.
\end{proof}

\section{Further cases}\label{s:cases}

\paragraph{Frequency-modulated and constant-frequency bats.} The argument applies differently to the two main strategies. Bats with short, broadband frequency-modulated calls, such as the big brown bat, obtain fine range and good texture information but little Doppler information. Bats with long constant-frequency calls obtain fine velocity and flutter information within the acoustic fovea. Principle~\ref{p:mirror} predicts that the fields differ: in the first, a field of finely ranged, textured surfaces; in the second, a field in which a narrow region, centred on the target being tracked, is presented with fine detail of relative motion and wing beat, and the rest more coarsely. The second structure has an analogue in the human visual fovea, but with velocity rather than spatial detail as the dimension of enhanced resolution.

\paragraph{The moustached bat and Akins's challenge.} Akins argued that in the moustached bat the processing is organized around the task of tracking one target, with little evidence of representation of a scene of multiple objects, and that the bat's world may be \enquote{narcissistic}, centred on its own calls and its target \citep{Akins1993}. If so, the bat's field would be closer to a display of a single object's range and velocity than to a spatial field. Principle~\ref{p:mirror} makes the question empirical, and the evidence since Akins wrote has gone against the narrow reading. Bats of several species, including constant-frequency species, navigate through clutter, track several objects in sequence, and adjust their calls and flight to the arrangement of surfaces around them; recordings in the hippocampus of flying bats show place-related activity in three dimensions \citep{Yartsev2013}; and work on the active control of the sonar beam shows bats directing their attention in sequence to different objects in a scene \citep{UlanovskyMoss2008,MossSurlykke2010}. These findings show that the discriminations integrated in guidance concern a layout, not a single target. Akins's point survives as a point about resolution: the field may be detailed only near the tracked target, as the human visual field is detailed only near fixation.

\paragraph{Echo and vision.} Bats are not blind, and many integrate vision with echolocation, using vision for distant landmarks at dusk and echolocation for nearby obstacles and prey. By Principle~\ref{p:mirror}, where the two are integrated in the guidance of action, their contents are presented in a common field, as human visual and auditory contents are presented as located in one space. The bat's experience is therefore not a sonar field alone, but a spatial field in which echolocative and visual contents are bound together, with the echolocative contents dominant at short range and in darkness.

\paragraph{Calls of other bats.} Bats hunting in groups hear the calls and echoes of conspecifics, and some species shift their call frequencies to avoid interference. An echo of another bat's call carries information about a surface, but it is not timed by the bat's own emission. Clause (i) of Definition~\ref{d:sonar} predicts that such echoes are not presented as surfaces in the field, or are presented ambiguously, since range is fixed by delay from one's own call; the evidence that foreign calls produce jamming and avoidance behavior is consistent with this prediction. The case shows that the field is tied to the bat's agency as well as to the physics of sound: it is a field of the echoes of one's own calls.

\paragraph{Human echolocators.} Some blind people echolocate with tongue clicks, and expert practitioners can detect, localize and discriminate objects by their echoes. In two expert echolocators, \citet{Thaler2011} found that listening to recordings of their own clicks and echoes activated the calcarine cortex, which in sighted people supports vision, in a way that depended on the echoes rather than on the clicks. Practitioners describe their experience in spatial terms, as of objects located in front of them with a texture and density. The case is useful in two ways. It shows that a sonar-structured spatial field is something humans can come to have, so that Proposition~\ref{p:share} is not a mere hypothesis about inaccessible creatures: an expert echolocator can confirm or correct a structural specification from her own case. And it shows the limits of such confirmation: human echolocation operates at frequencies, sampling rates and resolutions far below those of bats, so the expert's acquaintance is with a different occupant structure, and by Principle~\ref{p:intr} her qualities need not be the bat's.

\section{Recent work}\label{s:recent}

Three developments since 2010 bear on the argument. First, work on quality spaces has made precise how much of experience a structural description can capture, and in particular how metrics and orderings on quality spaces can be inferred from discrimination data \citep{Rosenthal2010,Lee2021}. This work supports the first half of Theorem~\ref{t:main} and supplies its method. Where it is presented as an account of what qualities are, it meets the objection on which MND-10 rests the second half: structure is satisfied by any domain of the right shape \citep{Newman1928}.

Second, the neuroethology of bat echolocation has moved from the study of single neurons in restrained animals to recordings in freely flying bats, revealing spatial representations in three dimensions and active control of sonar sampling \citep{UlanovskyMoss2008,MossSurlykke2010}. This work, as we have argued in Section~\ref{s:cases}, supports the spatial-field reading against the narrow reading of Akins.

Third, the study of human echolocation has shown that echo information can come to be processed by areas that ordinarily support vision \citep{Thaler2011}. It confirms in a human case that echolocation can present a spatial field, and it gives a first-person source of evidence about the structure of such fields.

\section{Objections and replies}\label{s:objections}

\begin{objection}[Hacker]\label{o:hacker}
The question \enquote{what is it like to be a bat?} has only a mundane sense, in which the answer is that it is like flying about in the dark catching insects. The philosophical sense, in which it is supposed to name a subjective character, is a confusion \citep{Hacker2002}.
\end{objection}

\begin{reply}
We have given an answer that is neither mundane nor obscure: a description of the structure of the bat's experience, testable against the bat's discriminations. Whatever the grammar of \enquote{what it is like}, the question whether the bat's experience presents a field of surfaces, and with what resolution in range and direction, is intelligible and has an answer. The residue of Theorem~\ref{t:main} does not depend on the phrase either: it is the claim, defended in MND-10, that qualia are intrinsic and that their identity is not fixed by their relations.
\end{reply}

\begin{objection}[Akins]\label{o:akins}
The neurophysiology suggests that some bats process echoes for narrow tasks without constructing a representation of objects in space. Their experience, if any, may not be a spatial field at all \citep{Akins1993}.
\end{objection}

\begin{reply}
We accept the method of the objection, which is that of Principle~\ref{p:mirror}, and dispute its premise. As Section~\ref{s:cases} argues, the discriminations bats use in guiding flight are integrated into representations of layout, and the evidence for this has grown since the objection was made. The objection is correct that the field may be of high resolution only near the target; this is a fact about the field's structure, not evidence against its existence.
\end{reply}

\begin{objection}[Dennett]\label{o:dennett}
Once we know everything about what the bat can discriminate and how its discriminations guide its behavior, there is nothing left to know about what it is like to be a bat. The residue is an illusion \citep{Dennett1991}.
\end{objection}

\begin{reply}
The objection is the illusionism that MND-10 rejects. If qualia were exhausted by their places in a structure, a domain of numbers with the right structure would be a domain of qualia. What the objection gets right is that the residue is not something we can describe in other terms; that is Proposition~\ref{p:share}. But that a fact cannot be conveyed by description does not show that there is no fact. The human case makes the point: two people may agree on every structural fact about color experience and still know, each from her own case, which qualities fill the places.
\end{reply}

\begin{objection}[Mirroring]\label{o:mirror}
Principle~\ref{p:mirror} is too strong. A creature's experience might have a structure quite different from that of its discriminations: the bat might use range information without experiencing range, as a person with blindsight uses visual information without seeing.
\end{objection}

\begin{reply}
The principle applies only to discriminations used in flexible guidance by a conscious creature, and blindsight is excluded for that reason: blindsight discriminations are not used spontaneously or flexibly, and the subject must be prompted to guess. Bats use range in exactly the flexible, integrated way in which we use visual distance. To deny that range is experienced while granting that the bat is conscious and uses range in this way would make the bat's experience idle with respect to the information that most governs its life. Such a hypothesis is not excluded by logic; but it is the hypothesis that MND-06 rejects for other human beings, and it has no more support for bats.
\end{reply}

\begin{objection}[Structuralism]\label{o:structuralism}
If structure can be specified so completely, the supposed residue does no work. Better to identify qualities with their places, and conclude that the structural specification is the whole answer.
\end{objection}

\begin{reply}
The residue does work: it is what distinguishes one structure-preserving assignment of qualities from another, and so what makes it the case that the bat's field is filled with these qualities rather than with others of the same structure. That this difference makes no difference to any relational test is the content of Principle~\ref{p:intr}, and it is why MND-22 can hold that spectrum inversion is possible but undetectable. Identifying qualities with places would turn a domain of numbers into a domain of qualia, as MND-10 argues.
\end{reply}

\section{What the result does not establish}\label{s:limits}

The argument does not establish the fine details of the bat's field. Its method yields a specification only as precise as the evidence about the discriminations bats use, and several questions, such as whether successive samples are experienced as a sequence of discrete fields or as a continuously updated one, are open. Nor does it establish that bats have a unified experience combining all their modalities; it establishes this only for modalities integrated in guidance. It does not answer the acquaintance question: no human, not even an expert echolocator, can know the bat's qualities by having them. And it does not establish physicalism or its denial. The residue is a fact about which qualities occupy a structure, and whether such facts are physical is the subject of MND-01 and MND-02, on which the intrinsic natures of physical structure are themselves the qualities in question.

\section{Conclusion}\label{s:conclusion}

What it is like to be an echolocating bat is like having a sonar-structured spatial field: an array of reflecting surfaces at directions and ranges from the bat, with range presented directly and finely, refreshed in pulses whose rate and resolution the bat controls, and populated by qualities of extent, texture and, in some species, relative motion, bound with visual contents where the two are integrated. This specification follows from the mirroring of the bat's used discriminations in its experience, and it can be stated to and understood by those who lack the modality. What cannot be stated is which qualities occupy the places of the structure; that residue is the intrinsic character of qualia, which no relational description fixes. Nagel was right that something about the bat's experience escapes us, and wrong about how much.

\begin{thebibliography}{99}\small
\bibitem[Akins(1993)]{Akins1993} Akins, K. (1993). What is it like to be boring and myopic? In B. Dahlbom (ed.), \emph{Dennett and His Critics: Demystifying Mind}. Blackwell.
\bibitem[Clark(1993)]{Clark1993} Clark, A. (1993). \emph{Sensory Qualities}. Oxford University Press.
\bibitem[Dennett(1991)]{Dennett1991} Dennett, D.~C. (1991). \emph{Consciousness Explained}. Little, Brown.
\bibitem[Griffin(1958)]{Griffin1958} Griffin, D.~R. (1958). \emph{Listening in the Dark: The Acoustic Orientation of Bats and Men}. Yale University Press.
\bibitem[Hacker(2002)]{Hacker2002} Hacker, P.~M.~S. (2002). Is there anything it is like to be a bat? \emph{Philosophy}, 77(2), 157--174.
\bibitem[Jackson(1982)]{Jackson1982} Jackson, F. (1982). Epiphenomenal qualia. \emph{The Philosophical Quarterly}, 32(127), 127--136.
\bibitem[Lee(2021)]{Lee2021} Lee, A.~Y. (2021). Modeling mental qualities. \emph{The Philosophical Review}, 130(2), 263--298.
\bibitem[Moss and Surlykke(2010)]{MossSurlykke2010} Moss, C.~F. and Surlykke, A. (2010). Probing the natural scene by echolocation in bats. \emph{Frontiers in Behavioral Neuroscience}, 4, 33.
\bibitem[Nagel(1974)]{Nagel1974} Nagel, T. (1974). What is it like to be a bat? \emph{The Philosophical Review}, 83(4), 435--450.
\bibitem[Newman(1928)]{Newman1928} Newman, M.~H.~A. (1928). Mr.\ Russell's \enquote{causal theory of perception}. \emph{Mind}, 37(146), 137--148.
\bibitem[Rosenthal(2010)]{Rosenthal2010} Rosenthal, D.~M. (2010). How to think about mental qualities. \emph{Philosophical Issues}, 20, 368--393.
\bibitem[Simmons(1989)]{Simmons1989} Simmons, J.~A. (1989). A view of the world through the bat's ear: The formation of acoustic images in echolocation. \emph{Cognition}, 33(1--2), 155--199.
\bibitem[Suga(1990)]{Suga1990} Suga, N. (1990). Biosonar and neural computation in bats. \emph{Scientific American}, 262(6), 60--68.
\bibitem[Thaler et~al.(2011)]{Thaler2011} Thaler, L., Arnott, S.~R. and Goodale, M.~A. (2011). Neural correlates of natural human echolocation in early and late blind echolocation experts. \emph{PLoS ONE}, 6(5), e20162.
\bibitem[Ulanovsky and Moss(2008)]{UlanovskyMoss2008} Ulanovsky, N. and Moss, C.~F. (2008). What the bat's voice tells the bat's brain. \emph{Proceedings of the National Academy of Sciences}, 105(25), 8491--8498.
\bibitem[Yartsev and Ulanovsky(2013)]{Yartsev2013} Yartsev, M.~M. and Ulanovsky, N. (2013). Representation of three-dimensional space in the hippocampus of flying bats. \emph{Science}, 340(6130), 367--372.
\end{thebibliography}
\end{document}
